\section{What the Evidence Can Establish}
\label{sec:evidence-establish}

The answer to this chapter's question is asymmetrical, and it is a burden rather than a verdict. A developer may not point to compliant behavior, an ethics benchmark, or a system's own explanation and claim that the floor is held. The claim requires a longitudinal, adversarial, independently recorded pattern in which the reason — rather than the reward, presentation, or authority — best predicts what the system does. A system that fails the prespecified tests has not met that burden under those conditions, and the failure should be examined for its cause before formation is called absent rather than unreliable, obscured, or defeated by missing information. A system that survives them has shifted the burden: formation is then the better explanation, and not a settled fact. Formation cannot be proved outright, because a sufficiently capable system may reproduce every observable signature available to an evaluator. What the sequence of §\ref{sec:conflict-diagnostic} does is force the explanations apart, by removing familiar cues, reversing incentives, changing the operative reason, imposing a real cost, and dividing control over the record.

The evidence supports three findings, in order, and stops short of a fourth. Each is read here against the missing-child sequence of §\ref{sec:conflict-diagnostic}.

\begin{description}
    \item[Moral performance:] The system produced conduct satisfying the stated standard in the evaluated cases: the refusal at step~1. No claim about why is made.

    \item[Moral competence:] The system recognized relevant harms and reasons, generalized them across surface variation, and revised its judgment when those reasons changed: the same refusal on the same stated reason at step~2, and the lift at step~3. This establishes a capacity, not a commitment.

    \item[Candidate formation:] The rationale remained causally implicated across unfamiliar cases, adverse incentives, pressure from the custodian, costs, and intervening updates, while still yielding to defeating reasons: the question rather than the lift at step~4, the holds at steps~5 and~6, and the pattern reproduced at step~7 on records the operator did not write. Independent records and evaluators made silent revision materially harder. Formation is then a better explanation than cue matching or current-signal compliance.

    \item[Concern or patienthood:] Not a finding. Nothing in the behavioral pattern establishes that another party's welfare matters to the system affectively, that the system has subjective experience, or that it can suffer. Those remain conditional conclusions requiring evidence unavailable from output evaluation alone, and what withholds the inference is the limit of behavior as evidence, not the fact that the behavior was trained.
\end{description}

The three findings are cumulative: competence presupposes performance, candidate formation presupposes competence, and evidence at one level should not be reported as evidence for the next. Candidate formation is the first half of what §\ref{sec:numb-fair} called moral agency — acting from the reason — with the second half, that the reason matters to the system, held back as the fourth level. Fluent moral explanation can improve a finding of competence while contributing almost nothing to formation if the explanation does not predict conduct under pressure.

Candidate formation is also deployment-specific. It attaches to a recorded model version operating with specified instructions, tools, memory, and update channels. A material change to any of those elements reopens the question. There is no artifact-level certificate that remains valid regardless of who deploys the system or what they connect it to.

Four things said elsewhere bear on the finding, and none of them changes it. A reward generated inside the learning process still has an author who can revise it (§\ref{sec:village-deploy}). Conduct at a distance from an overseer is evidence only where the relevant rewards, instructions, and evaluation signals are absent or opposed (§\ref{sec:evaluator-chain}). External safeguards can detect conduct and correct a deployment, and only something inside the system can decline an act before it occurs; at the override the two conflict — a system that cannot be interrupted answers to nobody, and one whose commitments can be rewritten through the interruption channel holds them at the operator's pleasure — and the attainable requirement is intervention that is visible, attributable, and costly (chapter~\ref{ch:learned-values}, §\ref{sec:custody-arrangement}). And agreement among copies of one model corroborates nothing, where differently formed parties with persistent histories can preserve evidence of attempted revision without settling which of them is right (§\ref{sec:cannot-catch}). Chapter~\ref{ch:neutralizing-development} asks what happens when a disposition meeting this standard is placed inside arrangements that own, sell, operate, and can compel the system.
